%! TEX root = thesis.tex

\chapter*{Abstract}
\begin{otherlanguage*}{american}

The pretraining of large language models (LLMs) consumes millions of GPU-hours
and emits hundreds of tonnes of CO$_2$ per run.  Since grid carbon intensity
varies strongly over time and region, training can be paused during dirty
periods and resumed when the grid is cleaner.  We built a simulation framework
and an adaptive optimizer that searches over pause thresholds, resume
thresholds, and start dates to find the best configurations.  Using historical
electricity data at five-minute resolution, we evaluated two frontier
models---DeepSeek~V3 and Kimi~K2---across five regions with different grid
characteristics.  Our main finding is that the gap between pause and resume
threshold should be small (at most 16~gCO$_2$eq/kWh): a wide gap never pays
off.  In highly variable grids like Germany and Italy, savings reached 43\%
and 33\%, but came with up to 200\% longer training time.  In stable grids
like China and the US, savings stayed below 10\% regardless of the chosen
policy.  All code is implemented as a browser-based application with an
interactive simulator and optimizer, available as open source.

\end{otherlanguage*}

\chapter*{Kurzfassung}
\begin{otherlanguage*}{ngerman}

Das Pretraining gro{\ss}er Sprachmodelle (LLMs) verbraucht Millionen von
GPU-Stunden und emittiert hunderte Tonnen CO$_2$ pro Durchlauf.  Da die
CO$_2$-Intensit\"at des Stromnetzes stark \"uber Zeit und Region variiert,
l\"asst sich das Training in Zeiten hoher Emissionen pausieren und bei
saubererem Strom fortsetzen.  Wir haben eine Simulationsplattform und einen
adaptiven Optimierer entwickelt, der Pausenschwellen,
Fortsetzungsschwellen und Starttermine durchsucht, um die besten
Konfigurationen zu finden.  Anhand historischer Stromdaten in
F\"unf-Minuten-Aufl\"osung haben wir zwei Spitzenmodelle---DeepSeek~V3 und
Kimi~K2---in f\"unf Regionen mit unterschiedlichen Netzeigenschaften
evaluiert.  Unser wichtigstes Ergebnis: Der Abstand zwischen Pausen- und
Fortsetzungsschwelle sollte klein sein (maximal 16~gCO$_2$eq/kWh); ein gro{\ss}er
Abstand zahlt sich nie aus.  In stark variablen Netzen wie Deutschland und
Italien erreichten wir Einsparungen von 43\% bzw. 33\%, allerdings bei bis zu
200\% l\"angerer Trainingsdauer.  In stabilen Netzen wie China und den USA
blieben die Einsparungen unabh\"angig von der gew\"ahlten Strategie unter
10\%.  Der gesamte Code ist als browserbasierte Anwendung mit interaktivem
Simulator und Optimierer implementiert und als Open Source verf\"ugbar.

\end{otherlanguage*}
